\section{The several Picard functors}
\source{fga-explained:page-0247}\source{fga-explained:page-0248}\source{fga-explained:page-0249}\source{fga-explained:page-0250}\source{fga-explained:page-0251}\source{fga-explained:page-0252}\source{fga-explained:page-0253}\source{fga-explained:page-0254}\source{fga-explained:page-0255}\source{fga-explained:page-0256}\source{fga-explained:page-0257}\source{fga-explained:page-0258}\source{fga-explained:page-0259}\source{fga-explained:page-0260}\source{fga-explained:page-0261}\source{fga-explained:page-0262}\source{fga-explained:page-0263}\source{fga-explained:page-0264}\source{fga-explained:page-0265}\source{fga-explained:page-0266}\source{fga-explained:page-0267}\source{fga-explained:page-0268}\source{fga-explained:page-0269}\source{fga-explained:page-0270}\source{fga-explained:page-0271}\source{fga-explained:page-0272}\source{fga-explained:page-0273}\source{fga-explained:page-0274}\source{fga-explained:page-0275}\source{fga-explained:page-0276}\source{fga-explained:page-0277}\source{fga-explained:page-0278}\source{fga-explained:page-0279}\source{fga-explained:page-0280}\source{fga-explained:page-0281}\source{fga-explained:page-0282}\source{fga-explained:page-0283}\source{fga-explained:page-0284}\source{fga-explained:page-0285}\source{fga-explained:page-0286}\source{fga-explained:page-0287}\source{fga-explained:page-0288}\source{fga-explained:page-0289}\source{fga-explained:page-0290}\source{fga-explained:page-0291}\source{fga-explained:page-0292}\source{fga-explained:page-0293}\source{fga-explained:page-0294}\source{fga-explained:page-0295}\source{fga-explained:page-0296}\source{fga-explained:page-0297}\source{fga-explained:page-0298}\source{fga-explained:page-0299}\source{fga-explained:page-0300}\source{fga-explained:page-0301}\source{fga-explained:page-0302}\source{fga-explained:page-0303}\source{fga-explained:page-0304}\source{fga-explained:page-0305}\source{fga-explained:page-0306}\source{fga-explained:page-0307}\source{fga-explained:page-0308}\source{fga-explained:page-0309}\source{fga-explained:page-0310}

\begin{definition}[absolute Picard functor]\label{fga-def-9-2-1}\dcref{9.2.1}\source{fga-explained:page-0262}\uses{}
For $f:X\to S$, the absolute Picard functor is
$\Pic_X(T)=\Pic(X_T)$ for locally noetherian $S$-schemes $T$.
\end{definition}\begin{definition}[relative Picard functor]\label{fga-def-9-2-2}\dcref{9.2.2}\source{fga-explained:page-0262}\uses{fga-def-9-2-1}
The relative Picard presheaf is $\Pic_{X/S}(T)=\Pic(X_T)/\Pic(T)$; write
$\Pic_{(X/S)}^{\mathrm{zar}}$, $\Pic_{(X/S)}^{\mathrm{et}}$ and
$\Pic_{(X/S)}^{\mathrm{fppf}}$ for its associated sheaves.
\end{definition}
\begin{theorem}[comparison]\label{fga-thm-9-2-5}\dcref{9.2.5}\source{fga-explained:page-0263}\uses{fga-def-9-2-2,fga-thm-4-23}
Assume $\OO_S\to f_*\OO_X$ is universally an isomorphism.  Then
\[
\Pic_{X/S}\hookrightarrow\Pic_{(X/S)}^{\mathrm{zar}}
\hookrightarrow\Pic_{(X/S)}^{\mathrm{et}}
\hookrightarrow\Pic_{(X/S)}^{\mathrm{fppf}}.
\]
If $f$ has a section, all three arrows are isomorphisms; a Zariski-local
section suffices for the last two, and an etale-local section suffices for the
last arrow.
\end{theorem}
\begin{proof}
An invertible sheaf trivialized after an fppf cover descends to a pullback from
the base by faithful flatness and the universal $f_*\OO_X=\OO_S$ condition,
proving injectivity.  With a section, rigidify along the section; automorphisms
of a rigidified line bundle are trivial, so fppf descent is effective.  The
local-section variants follow by sheaf locality.
\end{proof}
\begin{lemma}\label{fga-lem-9-2-7}\dcref{9.2.7}\source{fga-explained:page-0264}\uses{fga-thm-9-2-5}
If $\OO_S\simeq f_*\OO_X$, pullback $\Pic(S)\to\Pic(X)$ is fully faithful on
finite-rank locally free sheaves; its essential image consists of sheaves
$\mathcal M$ with $f^*f_*\mathcal M\simeq\mathcal M$.
\end{lemma}
\begin{definition}[rigidification]\label{fga-def-9-2-8}\dcref{9.2.8}\source{fga-explained:page-0265}\uses{fga-def-9-2-2}
When $f$ has a section $g:S\to X$, a $g$-rigidification of an invertible sheaf
$\mathcal L$ on $X_T$ is an isomorphism $\OO_T\simeq g_T^*\mathcal L$.
\end{definition}
\begin{lemma}\label{fga-lem-9-2-9}\dcref{9.2.9}\source{fga-explained:page-0265}\uses{fga-def-9-2-8}
Rigidified invertible sheaves modulo rigidified isomorphism map isomorphically
to $\Pic_{X/S}(T)$.
\end{lemma}
\begin{lemma}\label{fga-lem-9-2-10}\dcref{9.2.10}\source{fga-explained:page-0265}\uses{fga-def-9-2-8}
Under the universal $\OO_S\simeq f_*\OO_X$ hypothesis, every automorphism of a
rigidified invertible sheaf is the identity.
\end{lemma}

\section{Relative effective divisors}
\begin{definition}[relative effective Cartier divisor]\label{fga-def-9-3-3}\dcref{9.3.3}\source{fga-explained:page-0268}\uses{}
A relative effective divisor $D\subseteq X_T$ is an effective Cartier divisor
whose defining equation is a non-zero-divisor on every fiber of $X_T\to T$.
\end{definition}\begin{lemma}\label{fga-lem-9-3-4}\dcref{9.3.4}\source{fga-explained:page-0268}\uses{fga-def-9-3-3}
For a closed subscheme $D\subset X_T$, fiberwise regularity of a local defining
equation is equivalent to $\OO_D$ being flat over $T$; in that case $D$ is a
relative effective divisor.
\end{lemma}
\begin{definition}[divisor functor]\label{fga-def-9-3-6}\dcref{9.3.6}\source{fga-explained:page-0269}\uses{fga-def-9-3-3}
$\Div_{X/S}(T)$ is the set of relative effective divisors on $X_T$.
\end{definition}
\begin{theorem}\label{fga-thm-9-3-7}\dcref{9.3.7}\source{fga-explained:page-0269}\uses{fga-def-9-3-6,fga-thm-5-25,fga-lem-9-3-4}
If $X\to S$ is projective and flat, then $\Div_{X/S}$ is represented by an
open subscheme of $\Hilb_{X/S}$.
\end{theorem}
\begin{proof}
A divisor is a flat closed subscheme with invertible ideal.  In the universal
Hilbert family, invertibility of the ideal is open by finite presentation and
the flatness criterion; the open locus has exactly the divisor functor's
points.
\end{proof}
\begin{definition}[linear system]\label{fga-def-9-3-12}\dcref{9.3.12}\source{fga-explained:page-0271}\uses{fga-def-9-3-3}
For an invertible sheaf $\mathcal L$ on $X$, $\LinSys_{\mathcal L/S}$ is the
functor of pairs $(D,\mathcal M)$ with $D$ relative effective and
$\OO(D)\simeq\mathcal L\otimes f^*\mathcal M$.
\end{definition}
\begin{theorem}\label{fga-thm-9-3-13}\dcref{9.3.13}\source{fga-explained:page-0271}\uses{fga-def-9-3-12,fga-thm-5-7,fga-lem-9-3-4}
If $X/S$ is proper and flat with geometrically integral fibers, then
$\LinSys_{\mathcal L/S}$ is represented by the projective bundle
$\PP(f_*\mathcal L)$ on the open locus where cohomology and base change make
$f_*\mathcal L$ locally free; its universal point is the universal divisor.
\end{theorem}
\begin{proof}
A section of $\mathcal L$ up to multiplication by a unit defines an effective
divisor precisely when it is fiberwise a non-zero-divisor.  The projective
bundle of one-dimensional quotients of $f_*\mathcal L$ represents such sections;
the regularity lemma gives the open divisor locus and the universal property.
\end{proof}

\section{The Picard scheme}
\begin{definition}\label{fga-def-9-4-1}\dcref{9.4.1}\source{fga-explained:page-0272}\uses{fga-def-9-2-2}
If one of the four relative Picard functors is represented by a scheme, that
scheme is the Picard scheme $\Pic_{X/S}$.
\end{definition}
\begin{definition}[Abel map]\label{fga-def-9-4-6}\dcref{9.4.6}\source{fga-explained:page-0273}\uses{fga-def-9-3-6,fga-def-9-4-1}
The Abel map is the natural transformation
$A_{X/S}:\Div_{X/S}\to\Pic_{X/S}$, $D\mapsto\OO_{X_T}(D)$, whenever the
target is represented.
\end{definition}
\begin{theorem}[Picard scheme existence]\label{fga-thm-9-4-8}\dcref{9.4.8}\source{fga-explained:page-0273}\uses{fga-thm-9-2-5,fga-thm-9-3-7,fga-thm-9-3-13,fga-thm-5-25}
Assume $X\to S$ is flat, projective Zariski-locally on $S$, with integral
geometric fibers.  Then $\Pic_{X/S}$ exists, is separated and locally of finite
type over $S$, and represents $\Pic_{(X/S)}^{\mathrm{et}}$.  If $S$ is
noetherian and $X/S$ is projective, it is a disjoint union of open subschemes,
each an increasing union of open quasi-projective $S$-schemes.
\end{theorem}
\begin{proof}
Stratify the etale Picard sheaf by Hilbert polynomial and by a uniform
vanishing twist.  Each stratum is covered by the projective scheme
$\Div_{X/S}\times\Pic$ via the Abel map.  The relation on two divisors with
isomorphic associated line bundles is flat and proper; its quotient is a
quasi-projective scheme.  Increasing the twist gives the stated union, and
the quotient relation proves separatedness and the universal property.
\end{proof}
\begin{lemma}\label{fga-lem-9-4-9}\dcref{9.4.9}\source{fga-explained:page-0276}\uses{fga-thm-4-33}
If $Z\to P$ is an etale-sheaf surjection, $Z$ is quasi-projective, and
$R=Z\times_PZ\to Z$ is smooth and proper, then $P$ is represented by a
quasi-projective scheme and $Z\to P$ is smooth.
\end{lemma}
\begin{proposition}\label{fga-prop-9-4-17}\dcref{9.4.17}\source{fga-explained:page-0279}\uses{fga-thm-9-4-8}
Whenever $\Pic_{X/S}$ represents the fppf Picard sheaf, it is locally of finite
type over $S$.
\end{proposition}
\begin{theorem}\label{fga-thm-9-4-18-1}\dcref{9.4.18.1}\source{fga-explained:page-0279}\uses{fga-thm-9-4-8}
If $X/S$ is projective and flat with geometrically reduced connected fibers and
geometrically irreducible components of ordinary fibers, then $\Pic_{X/S}$
exists.
\end{theorem}
\begin{theorem}\label{fga-thm-9-4-18-2}\dcref{9.4.18.2}\source{fga-explained:page-0280}\uses{fga-thm-9-4-8}
If $X\to S$ is proper and $S$ integral, there is a nonempty open $V\subset S$
such that $\Pic_{X_V/V}$ exists and is a disjoint union of open
quasi-projective subschemes.
\end{theorem}
\begin{corollary}\label{fga-cor-9-4-18-3}\dcref{9.4.18.3}\source{fga-explained:page-0280}\uses{fga-thm-9-4-18-2}
For a complete scheme $X$ over a field $k$, $\Pic_{X/k}$ exists, represents the
fppf Picard sheaf, and is a disjoint union of open quasi-projective schemes.
\end{corollary}
\begin{theorem}\label{fga-thm-9-4-18-4}\dcref{9.4.18.4}\source{fga-explained:page-0281}\uses{fga-thm-5-22}
For a proper $X/S$ and coherent $\mathcal T$, the subfunctor of $S$-schemes on
which $\mathcal T$ is flat is represented by an unramified finite-type
$S$-scheme.
\end{theorem}
\begin{theorem}\label{fga-thm-9-4-18-5}\dcref{9.4.18.5}\source{fga-explained:page-0281}\uses{fga-thm-9-4-18-4}
For a surjective map of proper $S$-schemes with $S$ integral, after shrinking
$S$ the induced map on fppf Picard sheaves is representable by quasi-affine
maps of finite type.
\end{theorem}
\begin{theorem}\label{fga-thm-9-4-18-6}\dcref{9.4.18.6}\source{fga-explained:page-0282}\uses{fga-thm-9-2-5,fga-thm-9-4-18-4}
If $X\to S$ is flat, proper and finitely presented as an algebraic space and
$\OO_S\to f_*\OO_X$ is universally an isomorphism, the fppf Picard functor is
an algebraic space locally of finite presentation over $S$.
\end{theorem}

\section{The connected component of the identity}
\begin{lemma}\label{fga-lem-9-5-1}\dcref{9.5.1}\source{fga-explained:page-0285}\uses{fga-def-9-4-1}
For a group scheme $G$ locally of finite type over a field, the identity
component $G^0$ is open and closed, geometrically irreducible, a subgroup
scheme, and formation of $G^0$ commutes with extension of the ground field.
If $G$ has a geometrically reduced open, then $G$ is smooth.
\end{lemma}
\begin{proposition}\label{fga-prop-9-5-3}\dcref{9.5.3}\source{fga-explained:page-0286}\uses{fga-lem-9-5-1,fga-thm-9-4-8}
For $X/k$ proper, the identity component $\Pic^0_{X/k}$ is a subgroup scheme
of $\Pic_{X/k}$ and is geometrically connected.
\end{proposition}
\begin{theorem}[Picard variety]\label{fga-thm-9-5-4}\dcref{9.5.4}\source{fga-explained:page-0286}\uses{fga-prop-9-5-3,fga-thm-9-3-13}
If $X/k$ is proper over a field and $\Pic_{X/k}$ exists, then $\Pic^0_{X/k}$
is a projective group scheme; over characteristic zero it is reduced and hence
an abelian variety.
\end{theorem}
\begin{corollary}\label{fga-cor-9-5-5}\dcref{9.5.5}\source{fga-explained:page-0287}\uses{fga-thm-9-5-4}
For a proper geometrically integral variety over an algebraically closed field,
$\Pic^0$ is the Jacobian/Picard variety and is an abelian variety in
characteristic zero.
\end{corollary}
\begin{theorem}\label{fga-thm-9-5-11}\dcref{9.5.11}\source{fga-explained:page-0291}\uses{fga-thm-9-3-13,fga-thm-9-4-8}
The Abel map from the relative divisor scheme to $\Pic_{X/k}$ is smooth at a
divisor exactly when the corresponding line bundle has the expected vanishing;
its image contains the identity component, and $\Pic^0$ is generated by
differences of effective divisors.
\end{theorem}
\begin{definition}[Picard components]\label{fga-def-9-5-9}\dcref{9.5.9}\source{fga-explained:page-0290}\uses{fga-def-9-4-1}
For a component $C$ of $\Pic_{X/k}$, its translate and its inverse are the
components obtained by tensor product and duality; the component containing the
identity is $\Pic^0$.
\end{definition}
\begin{proposition}\label{fga-prop-9-5-19}\dcref{9.5.19}\source{fga-explained:page-0295}\uses{fga-thm-9-2-5,fga-thm-6-1-19}
The tangent space of $\Pic_{X/k}$ at the identity is $H^1(X,\OO_X)$, and the
cotangent space is $H^0(X,\Omega_X^1)$.
\end{proposition}
\begin{proposition}\label{fga-prop-9-5-20}\dcref{9.5.20}\source{fga-explained:page-0296}\uses{fga-prop-9-5-19,fga-thm-9-5-4}
The reduced identity component is smooth and its dimension is
$h^1(X,\OO_X)$; in characteristic zero it equals $\Pic^0$.
\end{proposition}

\section{The torsion component of the identity}
\begin{definition}\label{fga-def-9-6-1}\dcref{9.6.1}\source{fga-explained:page-0301}\uses{}
For invertible sheaves $\mathcal L,\mathcal M$ on a projective variety over a
field, $\mathcal L$ is algebraically (respectively numerically) equivalent to
$\mathcal M$ when a nonzero tensor power is algebraically equivalent
(respectively when all intersection degrees on complete curves agree).
\end{definition}\begin{definition}\label{fga-def-9-6-2}\dcref{9.6.2}\source{fga-explained:page-0301}\uses{}
A family of invertible sheaves is bounded if its members occur as fibers of one
invertible sheaf on $X\times T$ for a finite-type scheme $T$.
\end{definition}\begin{theorem}\label{fga-thm-9-6-3}\dcref{9.6.3}\source{fga-explained:page-0301}\uses{fga-def-9-6-1,fga-def-9-6-2}
For projective $X$ over an algebraically closed field and invertible $\mathcal L$,
the following are equivalent: $\mathcal L$ is algebraically equivalent to
$\OO_X$; it is numerically equivalent to $\OO_X$; the powers $\{\mathcal L^n\}$
form a bounded family; $\chi(\mathcal F\otimes\mathcal L)=\chi(\mathcal F)$
for every coherent $\mathcal F$; the equality holds for every complete integral
curve; and every $\mathcal L^n\otimes\OO_X(1)$ is ample.  The additional
intersection-number formulations in the source are equivalent when $X$ is
irreducible.
\end{theorem}
\begin{proof}
Algebraic equivalence implies numerical equivalence and constancy of Euler
characteristic in connected families.  Numerical equivalence gives boundedness
by uniform regularity and the Quot scheme.  Conversely bounded powers force two
powers into one connected component, hence algebraic equivalence.  The curve
criterion is Riemann--Roch; ampleness for every twist forces zero degree on every
curve.  The intersection formulations follow from the Hilbert polynomial.
\end{proof}
\begin{lemma}\label{fga-lem-9-6-5}\dcref{9.6.5}\source{fga-explained:page-0303}\uses{fga-def-9-6-1}
For a coherent $\mathcal F$ on projective $X$, there is a polynomial bound on
$\dim H^i(X,\mathcal L\otimes\mathcal F(n))$ uniform over invertible sheaves
$\mathcal L$ numerically equivalent to $\OO_X$.
\end{lemma}
\begin{lemma}\label{fga-lem-9-6-6}\dcref{9.6.6}\source{fga-explained:page-0304}\uses{fga-lem-9-6-5,fga-thm-5-3}
There is an integer $m$ such that every invertible sheaf numerically equivalent
to $\OO_X$ is $m$-regular and has the same Hilbert polynomial as $\OO_X$.
\end{lemma}
\begin{definition}[torsion component]\label{fga-def-9-6-8}\dcref{9.6.8}\source{fga-explained:page-0305}\uses{fga-lem-9-5-1}
For a commutative group scheme $G/S$ and $n>0$, $G^{\tau}$ is the union of the
components of the fibers whose $n$-th power lies in the identity component for
some $n$; equivalently it is the union of the inverse images of $G^0$ under all
power maps.
\end{definition}
\begin{lemma}\label{fga-lem-9-6-9}\dcref{9.6.9}\source{fga-explained:page-0305}\uses{fga-def-9-6-8}
For a commutative group scheme locally of finite type, $G^\tau$ is an open group
subscheme and commutes with field extension; if it is quasi-compact, it is closed
and of finite type.
\end{lemma}
\begin{proposition}\label{fga-prop-9-6-12}\dcref{9.6.12}\source{fga-explained:page-0306}\uses{fga-lem-9-6-9,fga-thm-9-6-3,fga-thm-9-4-8}
If $\Pic_{X/k}$ exists and $X$ is projective, $\Pic^\tau_{X/k}$ is a closed
finite-type open subgroup scheme, compatible with extension of $k$.
\end{proposition}
\begin{theorem}\label{fga-thm-9-6-16}\dcref{9.6.16}\source{fga-explained:page-0307}\uses{fga-prop-9-6-12,fga-thm-9-4-8}
If $X\to S$ is projective Zariski-locally, flat, with integral geometric fibers,
then the torsion component $\Pic^\tau_{X/S}$ is of finite type over $S$.
\end{theorem}
\begin{corollary}\label{fga-cor-9-6-17}\dcref{9.6.17}\source{fga-explained:page-0308}\uses{fga-thm-9-6-16}
For projective flat families with integral geometric fibers,
$\Pic^\tau_{X/S}$ is open and closed and has finite-type connected components.
\end{corollary}
\begin{theorem}\label{fga-thm-9-6-20}\dcref{9.6.20}\source{fga-explained:page-0309}\uses{fga-thm-9-6-3,fga-lem-9-6-6,fga-thm-9-4-8}
If $X/S$ is projective and flat with integral geometric fibers, then for every
Hilbert polynomial $P$ the locus $\Pic^P_{X/S}$ of invertible sheaves with
fiberwise Hilbert polynomial $P$ is of finite type over $S$; in particular the
components of $\Pic^\tau$ are bounded.
\end{theorem}
\begin{corollary}\label{fga-cor-9-6-25}\dcref{9.6.25}\source{fga-explained:page-0310}\uses{fga-thm-9-6-20}
The connected components of $\Pic_{X/S}$ are open and closed and the union of
components with a fixed Hilbert polynomial is quasi-compact.
\end{corollary}
